\section{Flags, parabolics and coset coordinates}
All type sets are finite. Groups need not be finite unless this is stated.
A typed incidence system has a reflexive symmetric incidence relation,
with distinct elements of the same type never incident. A flag is a
pairwise incident set. A chamber has every type. A geometry is such a
system in which every flag extends to a chamber. It is firm if every
nonchamber flag has at least two chamber extensions. A residue consists of
the elements incident with a fixed flag and having a type absent from
that flag. Residual connectedness means that every residue of rank at
least two has a connected incidence graph. The specified group action
is part of the data.

For subgroups $H_i\leq G$, indexed by $I$, write
\[
 H_J=\bigcap_{j\in J}H_j,\qquad H_\varnothing=G,\qquad B=H_I.
\]
The elements of type $i$ are the left cosets of $H_i$, and two cosets of
different types are incident if they intersect. The action is left
multiplication. Inversion converts this convention to the right cosets
of the sources. We will use the following three conditions:
\begin{description}
\item[(G1)] The indexed intersections $H_J$ are all distinct.
\item[(G2)] If $|I\setminus J|\geq2$, then
$H_J=\gen{H_{J\cup\{k\}}:k\in I\setminus J}$.
\item[(G3)] Every family of cosets of different types with pairwise
nonempty intersections has a common element.
\end{description}

\noindent\textbf{Proposition UB.Coset (the complete coset interface).}
Conditions (G1)--(G3) hold if and only if the coset system is firm,
residually connected and flag-transitive under its given action.
The residue at the standard flag of type $J$ is the coset geometry of
$H_J$ with subgroups $H_{J\cup\{k\}}$, $k\notin J$.

\begin{proof}
Under (G3), a flag has a representative $g$ common to its cosets and is
therefore contained in the chamber $(gH_i:i\in I)$. Translating $g$ to
the identity shows flag transitivity. Elements of type $k$ incident with
the standard flag of type $J$ can, again by (G3), be represented by an
element of $H_J$. Two such representatives give the same coset precisely
when they differ by $H_J\cap H_k$. This proves the residue description.
Condition (G1) gives $H_{I\setminus\{i\}}>B$, so every panel has at least
two elements. Every shorter flag is contained in a panel; this proves
firmness.

For completeness, a coset incidence graph for subgroups $(L_i:i\in S)$
of a group $L$, with $|S|\geq2$, is connected exactly when the $L_i$
generate $L$. Every edge can be represented by a common group element.
Successive representatives along a path differ by elements of the
subgroup corresponding to their shared vertex. Thus a connected path
from the identity cosets to $gL_i$ places $g$ in $\gen{L_j:j\in S}$.
Conversely, for $a\in L_j$ and $i\ne j$, the vertices
$L_i,L_j=aL_j,aL_i$ give a path. Translating and concatenating these
paths for a word in the $L_j$ connects every coset to the standard ones;
the standard vertices are mutually incident. Applying this criterion
to the described residues proves residual connectedness from (G2).

Conversely, flag transitivity takes any pairwise incident family to the
standard flag of its types, so it has a common representative and (G3)
holds. Firmness and panel transitivity give an element
$s_i\in H_{I\setminus\{i\}}\setminus B$ for every $i$. Its intersection
with $H_i$ is $B$, hence $s_i\notin H_i$. If $J\ne K$, choosing a type
in one set difference shows that $s_i$ belongs to one of $H_J,H_K$ and
not to the other. This proves (G1). Residual connectedness, the residue
description and the graph criterion just proved give (G2).
\end{proof}

An action is primitive when it is transitive and has no nontrivial
invariant partition. A geometry with its group action is PRI if the
action on every type is primitive, and WPRI if the action on at least
one type is primitive. The residual versions, RPRI and RWPRI, apply
to each residue of positive rank, including the whole geometry. In coset
coordinates, RWPRI says
\[
\tag{G4}
J\subsetneq I\quad\Longrightarrow\quad
\text{some }H_{J\cup\{k\}}\text{ is maximal in }H_J,
\quad k\notin J.
\]
This is the elementary equivalence between primitivity on a coset space
and maximality of its stabilizer. It refers to the stabilizer's induced
action; a kernel of that action does not change its invariant partitions.

\noindent\textbf{Lemma UB.Flags (residual weak primitivity gives flags).}
A finite-rank geometry with a residually weakly primitive group action
is flag-transitive.
\begin{proof}
Fix a chamber. In the empty-flag residue choose a type on which the group
is primitive, hence transitive. Match the element of that type in any
second chamber to the fixed one. In the residue of the matched flag,
choose a remaining primitive type and use its flag stabilizer to match
the next element. At most $|I|$ steps match the chambers. Since every
flag extends to a chamber, transitivity follows for flags of every
fixed type. The typewise bijection from cosets of stabilizers to
elements preserves incidence: transitivity on pairs of incident
elements supplies a common representative.
\end{proof}

For $B\leq G$, a finite indexed family $(s_i:i\in I)$ is independent
relative to $B$ if
$s_i\notin\gen{B,s_j:j\ne i}$ for every $i$. If it also generates
$G$ together with $B$, it is a relative generating family. For finite
$G$ let $\mu(G,B)$ and $\mu'(G,B)$ be the largest cardinalities of
relative independent generating families and relative independent
families, respectively. Omit $B$ if $B=1$. Independence for a group
acting on another group means independence under generation by all
images of the other family members; it is denoted $\mu'_G(N)$ in the
normal-quotient application.
